\chapter{Linearization of the Euler-Voigt System}
\label{app:lin:euler}

In this Appendix we derive the linearization of the Euler-Voigt equations. It is used to obtain the adjoint system and the gradient of the objective functional.
\\~\\
\noindent The Euler-Voigt equations are defined as,
\begin{equation*}
\begin{cases}
      -\alpha^2\partial_t\Delta \ua + \partial_t \ua + (\ua\cdot\bn)\ua+\bn \pa = 0\\
      \bn\cdot \ua = 0\\
      \left.\ua\right\vert_{t=0} = \bea.\\
    \end{cases}
\end{equation*}
Substituting in $\ua=\ua+\upa$, $\bea=\bea+\bepa$, and $\pa=\pa+\ppa$,
$$\begin{cases}
      -\alpha^2\partial_t\Delta (\ua+\upa) + \partial_t (\ua+\upa) + ((\ua+\upa)\cdot\bn)(\ua+\upa)+\bn (\pa+\ppa) = 0\\
      \bn\cdot (\ua+\upa) = 0\\
      \left.(\ua+\upa)\right\vert_{t=0} = \bea+\bepa.\\
    \end{cases}$$
By expanding each term we obtain,
$$\begin{cases}
      -\alpha^2\partial_t\Delta \ua-\alpha^2\partial_t\Delta \upa + \partial_t \ua+\partial_t\upa + \ua\cdot\bn \ua + \\ \ua\cdot\bn \upa +\upa\cdot\bn \ua + \upa\cdot\bn \upa +\bn \pa+ \bn \ppa = 0\\
      \bn\cdot (\ua+\upa) = 0\\
      \left.(\ua+\upa)\right\vert_{t=0} = \bea+\bepa.\\
    \end{cases}$$
Substituting in the Euler-Voigt equations we simplify to
$$\begin{cases}
      -\alpha^2\partial_t\Delta \upa + \partial_t\upa + \ua\cdot\bn \upa +\upa\cdot\bn \ua + \upa\cdot\bn \upa + \bn \ppa = 0\\
      \bn\cdot \upa = 0\\
      \left.\upa\right\vert_{t=0} = \bepa.\\
    \end{cases}$$
Linearizing the system by removing non-linear terms, we derive the linearized Euler-Voigt system
$$\begin{cases}
      -\alpha^2\partial_t\Delta \upa + \partial_t\upa + \ua\cdot\bn \upa +\upa\cdot\bn \ua + \bn \ppa = 0\\
      \bn\cdot \upa = 0\\
      \left.\upa\right\vert_{t=0} = \bepa.\\
    \end{cases}$$
This can be rearranged to the form
\begin{equation*}
\begin{cases}
      \partial_t \upa+(I-\alpha^2\Delta)^{-1}(\ua\cdot\bn \upa+\upa\cdot\bn \ua+\bn \ppa)=0\\
      \bn\cdot \upa = 0\\
      \left.\upa\right\vert_{t=0} = \bepa.\\
    \end{cases}
\end{equation*}